\section{Experiments}
\label{sec:results}
Grids e1--e10 (Table~\ref{tab:grids}, Appendix) comprise 806 cells (about 58 CPU-hours); a cell has $4$ worlds $\times$ $100$ assignments ($150$ for the coverage grid; $2\times60$ and $2\times50$ for the noiseless grids), so most points
average 400 estimates. Error bars are $\pm1.96$ Monte-Carlo SE, in units of the direct effect $\tau=1$. Unless stated, the default scenario is SBM, $n=5000$, mean degree 20, Bernoulli(0.5), $\gamma=2$ ($\tau_G=3$).
Verdicts follow the decision rules of Appendix~\ref{app:hyp} (directions from the plan, thresholds fixed after the grids ran). Grids e11--e21 (816 cells, 1622 in all) were designed after the first results were seen (e11--e13 with decision rules, e14--e21 as estimates with uncertainty only). Between-world variance of OLS$(z,f)$ is small (intraclass correlation of the error: median $\le0.004$, maximum $0.035$; 20-world grid e18, Appendix~\ref{app:post2}); with four worlds, bias$\pm1.96\,$SE (larger of replicate- and world-level SE) covered the 20-world reference in $0.97$ of draws, the world-level SE alone in $0.85$.

\subsection{Accuracy versus the baselines (Q1)}
\label{sec:res-acc}
\begin{figure}[tbp]
\centering\includegraphics[width=0.55\linewidth]{fig1_strength.pdf}
\caption{Bias (top) and RMSE (bottom, log scale) against interference strength $\gamma$ (linear-in-means truth, Bernoulli(0.5)). The difference in means has bias exactly $-\gamma$; the regression estimators are unbiased, with a fixed variance price. Other graphs: Fig.~\ref{fig:strength-all}.}
\label{fig:strength}
\end{figure}
With linear spillover the neighbour-fraction regression is unbiased on all six graphs, whereas the difference in means has bias slope $-1.000$ to $-0.999$ per unit of $\gamma$ (Fig.~\ref{fig:strength}); it wins on RMSE at $\gamma=0$ on all 12 graph--design combinations, and the regression overtakes it from $\gamma=0.5$ (BA, ER, SBM, Stanford3) or $\gamma=1$ (WS, ego-Facebook) under Bernoulli assignment.
The price is variance: at $\gamma=0$, when both are unbiased, the SD of OLS$(z,f)$ is 5.5--18.6 times that of the difference in means (BA 5.5, ER 6.5, Stanford3 9.2, SBM 10.1, WS 12.3, ego-Facebook 18.6), and a correctly specified regression pays the same price (ratio 0.84--1.00 to OLS$(z,f)$): this is the cost of extrapolating to the corners, not of misspecification. Hierarchical-bootstrap intervals are narrow (BA 5.5 [4.8, 6.3], ego-Facebook 18.6 [16.9, 20.7]), so the range is variation over graphs. It is not a constant: on an SBM family with $n\in\{10^3,5\cdot10^3,2\cdot10^4\}$, mean degree 10--40 and noise $\sigma\in\{0.5,1,2\}$ (e17) the ratio is 4.0--26.5 and scales as $d^{1.0}\sigma^{-0.31}n^{0.01}$ ($R^2=0.99$), and the interference strength $\gamma^*$ at which the regression overtakes the difference in means scales as $n^{-0.49}d^{1.01}\sigma^{0.09}$ ($R^2=0.997$), from $0.11$ to $2.2$. In the default scenario with \emph{concave} spillover (Table~\ref{tab:est}) the difference in means has RMSE 2.00 (67\% of $\tau_G$), OLS$(z,f)$ 0.72 (24\%), the oracle 0.59 (20\%), and $\kappa$-Hájek 1.21; the neighbour-count regression has bias $+1.0$ (RMSE 1.05), and richer specifications trade bias for variance: the cubic fit removes it ($-0.05$) but has RMSE 1.33, the 11-column fit RMSE 3.01 (SD 2.95 against 0.49). Separate-arm fitting and penalisation do not change this on the 48 Bernoulli cells of e2 (e16, Appendix~\ref{app:post2}): the fit $(z,f,zf)$ has a median RMSE $1.03\times$ that of OLS$(z,f)$ and is better in 19\% of cells; ridge and lasso on the 11 columns have median ratios $6.3$ and $6.5$ and are never better; boosting tuned by cross-validation has median ratio $2.0$ at Bernoulli(0.5) and is better in 21\% of cells, mostly with sigmoid truth, since trees cannot extrapolate to $f\in\{0,1\}$.
\textbf{Exposure-mapping baselines (H6; post-freeze grid e12, Appendix~\ref{app:post}, Fig.~\ref{fig:hajekfig}).} Hájek estimators are mostly unestimable at the degrees above, so we added low-degree graphs where they are estimable: with a correct exposure mapping they are unbiased when both extreme exposures are common (Bernoulli(0.5) and clusters: 30 of 30 cells) but not at Bernoulli(0.8) (1 of 13); with a wrong mapping they are biased by the size of the mismatch; they beat OLS$(z,f)$ in RMSE in only about 30\% of cells. The pre-stated rules gave ``unclear'' for both parts (31 of 43 and 44 of 65 cells). Under Louvain-cluster designs at degree 20--98 (e19) the full-exposure Hájek estimator is nearly unbiased ($|\text{bias}|\le0.05$) but has RMSE $1.19$--$1.47$ times that of OLS$(z,f)$ on SBM and ego-Facebook, and is competitive only on Stanford3 with concave truth ($0.89$--$0.96$).
\begin{table}[tbp]
\caption{Estimators on the default scenario (SBM, Bernoulli(0.5), $\gamma=2$, $\tau_G=3$): bias (MC SE), SD, RMSE. Source: \texttt{results/summary/e2.csv}; generated by \texttt{figures/scripts/make\_tables.py}.}
\label{tab:est}
\centering\scriptsize\setlength\tabcolsep{3pt}
\input{tab2_estimators.tex}
\end{table}

\subsection{Where the bias comes from, and why it does not average out (Q2)}
\label{sec:res-theory}
\begin{figure}[tbp]
\centering\includegraphics[width=0.7\linewidth]{fig2_bias_theory.pdf}
\caption{(a) Noiseless simulated bias of OLS$(z,f)$ against the closed form of Eq.~(\ref{eq:bias}) in 144 cells (SBM, BA, Stanford3; concave/convex shapes, curvature 0.5--8, $p\in\{0.1,0.25,0.5,0.75\}$). (b,c) Bias against $p$: lines are the closed form, markers simulation.}
\label{fig:theory}
\end{figure}
\textbf{The bias is the extrapolation of the local slope (checks H1, H2).} With a correctly specified model and no noise, the bias is zero to machine precision (max $|\text{bias}|=2.4\times10^{-11}$ over 120 cells); in noisy cells only 0.6\% of 156 cells exceed 3 SE. For OLS$(z,f)$ under nonlinear one-hop spillover the exact expression (\ref{eq:bias}) reproduces all 144 noiseless cells (correlation $1.0000$, median $|\text{dev}|=0.0012$, max $0.041$, all within 3 SE; Fig.~\ref{fig:theory}, Table~\ref{tab:theory}) over a range of biases from $-2$ to $+5$ ($\tau_G=3$). The sign follows the curvature of $h$ at $p$: concave spillover at small $p$ is over-, at large $p$ under-estimated, and the sign flips with $p$. The large-degree shortcut $\gamma(h'(p)-1)$ is a much poorer description: it deviates by a median 0.09 and up to 3.3, on Stanford3 by 0.5--0.8 on average, because the low-degree units dominate the exposure distribution (the exact formula needs the degree distribution).
\textbf{It is not a sampling error (H3).} In SBM with concave spillover the bias of OLS$(z,f)$ is $+2.53$ at $n=10^3$ and $+2.49$ at $n=2\times10^4$ for $p=0.1$ ($-0.57\to-0.52$ at $p=0.5$; log--log slopes $-0.01$, $-0.03$), whereas the SD falls as $n^{-1/2}$ (Fig.~\ref{fig:bias-n}); more data tighten a wrong answer. The bias grows with the curvature (median Spearman 0.92 over 12 graph--design groups). The bias is proportional to the interference strength (post-freeze grid e11, Appendix~\ref{app:post}): 14 of 16 graph--shape--$p$ groups pass the rule (R$^2\ge0.99$, intercept within 3 SE; the two failures are ego-Facebook), with slopes per unit $\gamma$ for OLS$(z,f)$ from $-0.71$ to $+1.26$ set by shape and $p$.
\begin{figure}[tbp]
\centering\includegraphics[width=0.64\linewidth]{fig3_bias_vs_n.pdf}
\caption{Bias (symlog) and SD (log--log) against $n$, concave spillover, three designs. The oracle (green dashed) is unbiased; the bias of the misspecified regressions is flat in $n$ while the SD falls with slope $-1/2$ (dotted).}
\label{fig:bias-n}
\end{figure}

\textbf{Variance (H4).} On Erd\H{o}s--R\'enyi graphs under Bernoulli(0.5) the SD of OLS$(z,f)$ follows $\sqrt{d/n}$ with fitted exponents $-0.501\pm0.005$ in $n$ and $+0.495\pm0.006$ in $d$ over 25 cells (Fig.~\ref{fig:var}, Table~\ref{tab:scaling}); the correctly specified model and the model with a log-degree term have the same exponents, but the 11-column specification does not ($-0.61$, $+1.18$: its variance explodes with degree). The dependence on $p$ follows $1/\sqrt{p(1-p)}$ on SBM and Stanford3 (SD at $p=0.05$ relative to $p=0.5$: 2.32 and 2.21 against a predicted 2.29), but not on ego-Facebook (1.80; log--log slope $0.59$ against 1), whose hubs break the homogeneous-degree argument (decision rule: unclear). Under Louvain clusters the scaling differs ($n^{-0.55}$, $d^{+0.37}$). Raising the graph-smooth baseline component from 0 to 1 to 3 raises the SD of OLS$(z,f)$ from 0.16 to 0.49 to 1.37 at $\sigma=1$ (e5): neighbours' shared baselines are an additional variance source. The heuristic with the harmonic-mean degree $d_h$ (Sec.~\ref{sec:setup}) holds across graphs when the errors are independent (e21, $\gamma=0$, noise only: observed over predicted SD between $0.93$ and $1.04$ on all six graphs, while the arithmetic mean degree is off by a factor 2.8 between graphs), but not for the default outcome: with its graph-smooth and log-degree baselines the ratio ranges from $1.3$ (ER) to $5.1$ (ego-Facebook), and neither $d_h$ nor $d$ orders the six graphs (Spearman correlation $0.03$). The \emph{scaling} in $n$, $d$ and $p$ is thus supported for independent errors and homogeneous graphs; the \emph{level} on a given graph depends on the correlation of the baseline, which the formula does not contain.
\begin{figure}[tbp]
\centering\includegraphics[width=0.66\linewidth]{fig4_variance_scaling.pdf}
\caption{(a) SD against $n$ for each mean degree $d$ (ER, Bernoulli(0.5)); dotted: fitted $C\sqrt{d/n}$. (b) $\mathrm{SD}/\sqrt{d/n}$ is flat for OLS$(z,f)$, not for the 11-column fit. (c) SD relative to $p=0.5$ against $p$.}
\label{fig:var}
\end{figure}

\subsection{Design: when cluster randomisation helps (H5)}
\label{sec:res-design}
Louvain-cluster randomisation removes most of the difference-in-means bias in every matched cell ($|\text{bias}|$ smaller in 100\% of 36 matched resolutions $\times$ graph $\times$ shape cells), whereas random clusters of the same sizes leave it at the Bernoulli value ($-2.0$): clusters help through the graph, not through coarse assignment (Fig.~\ref{fig:cluster}, Appendix). For the regression estimator the picture depends on the graph. On SBM and ego-Facebook the Louvain design has lower RMSE (SBM, concave: bias $-0.07$ vs.\ $-0.48$, RMSE $0.37$ vs.\ $0.76$; ego-Facebook, linear: RMSE $0.43$ vs.\ $0.98$), but on the dense Stanford3 graph (mean degree 98) it is \emph{reversed} (RMSE $0.54$ vs.\ $0.26$, linear), and the ordering replicates on a second seed (Louvain worse in 10 of 12 matched cells in the first seed and 11 of 12 in the second; the 12-cell mean RMSEs flip, so the claim is about the cell-wise ordering). That good clusters leave little exposure variation to the regression holds for SBM and ego-Facebook (within-arm variance of $f$ 0.005 and 0.004--0.04, against 0.013 and 0.023 under Bernoulli) but \emph{not} for Stanford3, where it stays at the Bernoulli value (0.021--0.028 against 0.023); there the SD inflation goes with the few clusters and the cluster-level dependence of the treated share (Appendix~\ref{app:post2}). Overall Louvain beats random clusters in 72\% of 36 matched cells (decision rule $\ge75\%$: unclear).

\subsection{Inference (H7)}
\label{sec:res-cov}
\begin{figure}[tbp]
\centering\includegraphics[width=0.65\linewidth]{fig6_coverage.pdf}
\caption{Empirical coverage of nominal 95\% intervals for OLS$(z,f)$ (filled) and the correctly specified regression (open) by variance estimator (HC0; cluster-robust, CL, at analysis resolutions; graph-HAC at 1 and 2 hops), spillover shape and design.}
\label{fig:cov}
\end{figure}
Intervals that treat units as independent (HC0) undercover whenever neighbours share exposures. At the default analysis resolution, Bernoulli(0.5), and linear spillover, OLS$(z,f)$ has mean coverage 0.66 for HC0 (0.32--0.89 across the six graphs; lowest on ego-Facebook), 0.86 for cluster-robust errors, 0.79 for 1-hop graph-HAC and 0.91 for 2-hop graph-HAC (Fig.~\ref{fig:cov}). Under misspecification no interval works: with concave spillover the same estimator covers 0.43 (HC0), 0.65 (cluster), 0.56 (HAC-1) and 0.70 (HAC-2) on average, because the bias is not in the standard error; the correctly specified regression reaches 0.86 (cluster) and 0.92 (HAC-2), HC0 only 0.68. Binomial MC SEs of single coverage values are 0.01--0.02 (lower bounds). Cells with a single analysis cluster (BA and ER at analysis resolution 0.5) are degenerate and excluded; the 2-hop HAC of this grid is over five graphs.
A repeat on the same six graphs, with the Stanford3 2-hop HAC added, a Louvain-block bootstrap ($B=200$) and coverage of both the GATE and the expectation $E[\hat\tau]$ (which removes the bias; e15, mean over six graphs $\times$ linear/concave truth) gives, for GATE/$E[\hat\tau]$ under Bernoulli(0.5): HC0 $0.54/0.66$, 1-hop HAC $0.66/0.78$, cluster-robust $0.75/0.86$, block bootstrap $0.76/0.87$, 2-hop HAC $0.82/0.92$; and with Louvain clusters $0.36/0.37$, $0.56/0.57$, $0.80/0.82$, $0.83/0.85$, $0.73/0.74$ (Appendix~\ref{app:post2}). The gap between the targets is the bias; the bootstrap equals cluster-robust errors under Bernoulli assignment, is the best under cluster designs, and no interval reaches 95\% on average. We resampled analysis clusters; the bootstrap of \citet{hanugander2023} was not implemented.

\subsection{Failure map and ranking of factors (Q2, Q3)}
\label{sec:res-fail}
\begin{figure}[tbp]
\centering\includegraphics[width=0.72\linewidth]{fig7_failure_map_relrmse_sbm.pdf}
\caption{Relative RMSE (RMSE$/|\tau_G|$) for each estimator, spillover shape or mechanism and design on SBM (B.5/B.1: Bernoulli 0.5/0.1; Lv: Louvain clusters). Hatched: not estimable. Other graphs: Appendix~\ref{app:figs}.}
\label{fig:fail}
\end{figure}
\textbf{Shape and treatment probability.} Median $|\text{bias}|/\tau_G$ of OLS$(z,f)$ is 0.01 (linear), 0.14 (concave), 0.14 (convex) and 0.34 (sigmoid) over 24 cells each, and the squared bias exceeds half the MSE in 0, 33, 38 and 67\% of cells (Fig.~\ref{fig:fail}). At $p=0.1$ the regression family is much worse than at $0.5$ (on SBM, concave: relative RMSE 0.86 for OLS$(z,f)$, 3.8 for OLS quad and 27 for the 11-column set) and the bias changes sign (sigmoid: $+2.0$ at $p=0.5$, $-1.4$ to $-1.5$ at $p=0.1$). Over 96 graph--shape--design cells OLS$(z,f)$ beats the difference in means in 85\%, but a correctly specified model beats OLS$(z,f)$ in 85\%, while quadratic features (14\%) and gradient boosting (18\%) rarely do: flexibility costs variance and extrapolates badly (OLS quad bias $-11$ at $p=0.1$, $n=2\times10^4$).
\textbf{Other mechanisms.} With two-hop interference fitted by one-hop features OLS$(z,f)$ has bias $-0.4$ (WS) to $-1.0$ (SBM, BA, ER, Stanford3), yet \emph{adding} the two-hop exposure (correct model) has larger RMSE on every graph (SBM $2.75$ vs.\ $1.04$): the variance of the extra regressor exceeds the bias it removes. Near the singular end of equilibrium propagation ($\rho=0.9$; $\tau_G\approx10$, relative bias of OLS$(z,f)$ $-0.74$) no estimator is safe (mean RMSE: OLS$(z,f)$ 7.8, difference in means 8.9, boosting 8.7, 11-column 9.8); $\rho=0.3$ is benign.
\textbf{Graphs.} The predicted penalty for heavy-tailed degrees was not observed (H8): for linear truth OLS$(z,f)$ RMSE is 0.26 on BA and Stanford3, 0.27 ER, 0.47 SBM, 0.51 WS and 0.92 ego-Facebook, so the two heavy-tailed graphs are among the best, but graph family is confounded with $n$ and mean degree here. A matched test (e14: $n=5000$, mean degree 20, linear truth, $\gamma=2$) removes the confounding: the SD of OLS$(z,f)$ is 0.255 on BA ($m=10$), 0.248 on a configuration model with the same degree sequence, 0.276 on ER, 0.240 on a 20-regular graph and 0.144--0.175 on power laws (exponents 3.0, 2.5, 2.2; $d_h$ 3.6--4.5), so heavier tails lower it, in the direction of the $d_h$ heuristic, and there is no penalty; under concave truth they also lower the bias ($-0.12$ to $-0.14$ against $-0.45$ to $-0.52$).
\begin{figure}[tbp]
\centering\includegraphics[width=0.96\linewidth]{fig11_crossover_anova.pdf}
\caption{(a) Crossover strength $\gamma^*$ at which OLS$(z,f)$ overtakes the difference in means in RMSE (SBM, linear truth, Bernoulli(0.5); from the fitted difference-in-means bias slope and the OLS RMSE) against mean degree, lines by $n$, markers by $\sigma$; open markers: $\gamma^*>2$, extrapolated. (b) SD ratio SD(OLS$_f$)/SD(DM) at $\gamma=0$ on the same cells, with the six default graphs (95\% hierarchical-bootstrap intervals). Printed: exponents of the log--log fits over 27 cells ($\pm$ SE). (c) First-order and total-effect shares of the variance of $\log$ RMSE(OLS$_f$), SBM factorial of 144 cells (e20); whiskers: 95\% parametric bootstrap.}
\label{fig:cross}
\end{figure}
\textbf{Which factors matter (C3).} One-factor sweeps give max/min RMSE ratios of OLS$(z,f)$ of $4.6$ for sample size, $4.5$ for spillover shape, $3.9$ for mean degree, $3.6$ for $p$, $3.2$ for graph family, $2.0$ for design and $1.5$ for noise (Fig.~\ref{fig:rank}, Appendix~\ref{app:figs}), but these depend on the width of each range and are not a ranking. A full factorial on SBM (e20, 144 cells; $n$, $d$, $p$, $\sigma$, graph-smooth baseline, shape) decomposes the variance of log RMSE of OLS$(z,f)$ into first-order/total-effect shares: graph-smooth baseline $0.22/0.33$, spillover shape $0.20/0.36$, $p$ $0.16/0.20$, $n$ $0.13/0.17$, mean degree $0.09/0.12$, noise $0.01/0.04$ (Fig.~\ref{fig:cross}c). The outcome side (baseline structure and shape) thus matters most and noise least, in these ranges and on SBM; graph family and design were not crossed, and the one-factor ratios remain the only evidence for them.
